---
artifact_kind: l3_flow_notebook
generated_at: '2026-04-28T16:33:55+08:00'
stage: L3
topic_slug: nl-test-qsgw
---
\documentclass[11pt,a4paper]{article}
\usepackage[utf8]{inputenc}
\usepackage[T1]{fontenc}
\usepackage{amsmath,amssymb,amsfonts}
\usepackage{physics}
\usepackage{bm}
\usepackage[margin=2.5cm]{geometry}
\setlength{\parindent}{0pt}
\setlength{\parskip}{0.5em plus 0.2em minus 0.1em}
\usepackage{enumitem}
\usepackage{booktabs}
\usepackage{longtable}
\usepackage{array}
\usepackage{xcolor}
\usepackage[most]{tcolorbox}
\usepackage{hyperref}
\hypersetup{colorlinks=true,linkcolor=blue!60!black,urlcolor=blue!60!black}

\definecolor{passgreen}{RGB}{34,139,34}
\definecolor{failred}{RGB}{178,34,34}
\definecolor{warnorange}{RGB}{204,119,34}
\definecolor{infoblue}{RGB}{47,85,151}
\definecolor{lightbg}{RGB}{245,245,250}

\newcommand{\statuspass}{\textcolor{passgreen}{\textbf{PASS}}}
\newcommand{\statusfail}{\textcolor{failred}{\textbf{FAIL}}}
\newcommand{\statusactive}{\textcolor{warnorange}{\textbf{ACTIVE}}}
\newcommand{\statusdeferred}{\textcolor{gray}{\textbf{DEFERRED}}}

\newtcolorbox{resultbox}[1][]{colback=infoblue!4,colframe=infoblue!70,
  fonttitle=\bfseries,boxrule=0.6pt,arc=2pt,breakable,title={#1}}
\newtcolorbox{warningbox}[1][]{colback=warnorange!4,colframe=warnorange!70,
  fonttitle=\bfseries,boxrule=0.6pt,arc=2pt,breakable,title={#1}}
\newtcolorbox{valbox}[1][]{colback=passgreen!4,colframe=passgreen!70,
  fonttitle=\bfseries,boxrule=0.6pt,arc=2pt,breakable,title={#1}}

\title{\textbf{Research Flow Notebook}\\[0.4em]
  \large Systematic Improvements of qsGW for Correlated Materials
  \quad \texttt{formal\_theory lane}
}
\author{AITP Protocol v3}
\date{2026-04-28}

\begin{document}
\maketitle
\begin{abstract}
This notebook records the complete research trajectory: sources, conventions,
derivation attempts, validated results, and open questions.
Mode: \texttt{explore}.
\end{abstract}

\tableofcontents
\newpage
\section{Research Question}

\begin{resultbox}[Bounded Question]
How do vertex corrections beyond the GW approximation systematically modify the quasi-particle band gap predicted by qsGW in moderately correlated materials?
\end{resultbox}

\noindent
\textbf{Scope:} One-shot vertex corrections (Gamma != 1) to the self-energy within the GW framework. Does NOT consider: finite-temperature effects, spin-orbit coupling, magnetic ordering, or full self-consistent vertex corrections. \quad
\textbf{Target quantities:} Direct band gap at Gamma point (eV), quasi-particle weight Z, self-energy correction Delta\_Sigma

\section{Source Landscape}

\begin{longtable}{>{\raggedright}p{3cm} p{5cm} p{2.5cm} p{2.5cm}}
\toprule
\textbf{Source} & \textbf{Title} & \textbf{Fidelity} & \textbf{Relation} \\
\midrule
\endhead
gruneis2014 & Gruneis A, Kresse G, Hinuma Y, Oba F (2014) Ionization poten & arxiv\_preprint &  \\
hedin1965 & Hedin L (1965) New Method for Calculating the One-Particle G & arxiv\_preprint &  \\
kotani2007 & Kotani T, van Schilfgaarde M, Faleev SV (2007) Quasiparticle & arxiv\_preprint &  \\
\bottomrule
\end{longtable}

\subsection*{Source Basis Analysis}
\subsection*{Source Basis}

\subsubsection*{Core Sources}
\begin{itemize}
\item hedin1965: Hedin 1965 foundational GW paper
\item kotani2007: qsGW method
\item gruneis2014: Vertex corrections
\end{itemize}

\subsubsection*{Peripheral Sources}
\begin{itemize}
\item biermann2015: DMFT+GW combination (deferred for now)
\end{itemize}

\subsubsection*{Source Roles}
hedin1965: Foundational. kotani2007: Direct dependency. gruneis2014: Direct dependency.

\subsubsection*{Reading Depth}
Full read for all core sources.

\subsubsection*{Why Each Source Matters}
Each provides one piece of the qsGW+vertex puzzle.



\section{Conventions \& Notation}

\subsection*{Convention Snapshot}

\subsubsection*{Notation Choices}
Hedin 1965 conventions as canonical:
\begin{itemize}
\item G(1,2): single-particle Green function
\item W(1,2): dynamically screened Coulomb interaction
\item Sigma(1,2): electron self-energy Sigma = iGW
\item Gamma(1,2;3): irreducible vertex function
\end{itemize}

\subsubsection*{Unit Conventions}
Atomic units. eV for reporting.

\subsubsection*{Sign Conventions}
\begin{itemize}
\item Time-ordered Greens functions
\item Self-energy: Sigma = iGW (Hedin convention)
\item Fourier: int domega/(2pi)
\end{itemize}

\subsubsection*{Metric Or Coordinate Conventions}
Real space, plane-wave basis.

\subsubsection*{Categorized Assumptions}
\paragraph{Mathematical assumptions}
Translational invariance, single Slater determinant reference.
\paragraph{Physical assumptions}
T=0, no SOC, weak-to-moderate correlation.
\paragraph{Notational assumptions}
Spin sum implicit, 1+ = t1+eta.

\subsubsection*{Canonical Notation}
Hedin 1965 canonical. All GW literature derives from it.

\subsubsection*{Unresolved Tensions}
Time-ordered vs retarded G for spectral properties.

\subsubsection*{L3 Discoveries}
(To be appended during derivation)



\noindent\textbf{Notation:} Hedin 1965: G=Green function, W=screened Coulomb, Sigma=self-energy, Gamma=vertex, P=polarizability

\noindent\textbf{Units:} Atomic units (Hartree/Bohr). eV for band gap reporting.

\section{Mode \& Session History}

\noindent
\begin{tabular}{ll}
\toprule
\textbf{Axis} & \textbf{Value} \\
\midrule
Topic slug & \texttt{nl-test-qsgw} \\
Status & \texttt{new} \\
Mode & \texttt{explore} \\
Lane & \texttt{formal\_theory} \\
L3 mode & \texttt{research} \\
L3 subplane & \texttt{derive} \\
\bottomrule
\end{tabular}

\subsection*{Session Log}
\subsection*{Session Log: Systematic Improvements of qsGW for Correlated Materials}

\subsubsection*{Sessions}



\section{Derivation Journey}

\noindent (No structured derivation recorded.)

\section{Synthesis \& Claims}

\section{Validation}

\begin{valbox}[cand-qsgw-vertex \quad partial\_pass]
\end{valbox}

\begin{valbox}[cand-qsgw-vertex\_v1 \quad partial\_pass]
\end{valbox}

\section{Negative Results \& Open Questions}

\noindent (No negative results or open questions formally recorded.)


\end{document}
